\chapter[Skalarne prawa zachowania]{Skalarne prawa zachowania: fale uderzeniowe, entropia i twierdzenie Krużkowa}
\label{chap:pde24}

\section{Od charakterystyk do rozwiązań słabych}
Niech $T>0$, $f\in C^2(\R)$ i $u_0\colon\R\to\R$ będzie
funkcją mierzalną. Rozważamy zagadnienie Cauchy'ego
\begin{equation}\label{eq:pde24-cl}
 u_t+\partial_x f(u)=0\quad\text{w }(0,T)\times\R,
 \qquad u(0,x)=u_0(x).
\end{equation}
Funkcję $f$ nazywamy \emph{strumieniem}. Dla gładkiego rozwiązania
równanie ma postać $u_t+f'(u)u_x=0$: wartości $u$ są stałe
wzdłuż charakterystyk o prędkości $f'(u)$. Po przecięciu
charakterystyk nie można na ogół kontynuować rozwiązania
w klasie $C^1$; musimy dopuścić nieciągłości.

Poniższe sformułowanie całkowe jest szczególnym
przypadkiem definicji~\ref{def:pde2-weak}.
Przyjmujemy dodatkowo lokalną ograniczoność rozwiązania,
odpowiednią dla teorii entropijnej.

\begin{uwaga}[Potencjał a prawo zachowania]
\label{rem:pde-v43-hj-entropy-bridge}
Dla $f(w)=w^2/2$ związek z równaniem
Hamiltona--Jacobiego opisano w
rozdziale~\ref{chap:pde23}. Pełna rachunkowa zgodność
\emph{lepkich} równań jest udowodniona w
Stwierdzeniu~\ref{prop:pde-v43-hj-viscous}.
Selekcja rozwiązania entropijnego po zaniku lepkości
wymaga dodatkowej zwartości i jest odrębnym etapem,
którego nie zastępuje formalne różniczkowanie
rozwinięcia charakterystyk.
\end{uwaga}

\begin{definicja}[Rozwiązanie słabe]
Niech $u_0\in L^1_{\mathrm{loc}}(\R)$ i
$u\in L^\infty_{\mathrm{loc}}([0,T)\times\R)$.
Funkcja $u$ jest słabym rozwiązaniem~\eqref{eq:pde24-cl}, jeżeli
\begin{equation}\label{eq:pde24-weak}
 \int_0^T\!\int_\R
  \bigl(u\varphi_t+f(u)\varphi_x\bigr)\,\mathrm d\lambda_1(x)\,\mathrm d\lambda_1(t)
 +\int_\R u_0(x)\varphi(0,x)\,\mathrm d\lambda_1(x)=0
\end{equation}
dla każdej $\varphi\in C_c^\infty([0,T)\times\R)$.
W globalnej teorii kontrakcji będziemy dodatkowo wymagać
$u\in C([0,T];L^1(\R))\cap L^\infty((0,T)\times\R)$.
W teorii lokalnej zastąpimy globalną ciągłość w $L^1$
ciągłością w $L^1_{\mathrm{loc}}$.
\end{definicja}

\begin{stwierdzenie}[Bilans na ruchomym przedziale]
\label{prop:pde24-balance}
Niech $u\in C^1$ spełnia~\eqref{eq:pde24-cl} i niech
$a,b\in C^1([0,T])$ oraz $a(t)<b(t)$. Wówczas
\[
 \frac{\mathrm d}{\mathrm dt}\int_{a(t)}^{b(t)}u(t,x)\,\mathrm d\lambda_1(x)
 =f(u(t,a(t)))-f(u(t,b(t)))
 +b'(t)u(t,b(t))-a'(t)u(t,a(t)).
\]
\end{stwierdzenie}
\begin{proof}
Reguła Leibniza daje dwa składniki pochodzące z przesuwających
się końców przedziału. W pozostałej całce podstawiamy
$u_t=-\partial_xf(u)$ i stosujemy twierdzenie Newtona--Leibniza.
\end{proof}

\begin{definicja}[Wariacja całkowita i przestrzeń $BV$]
Niech $w\in L^1(\R)$ będzie funkcją o wartościach
rzeczywistych. Definiujemy wariację całkowitą
\[
 \operatorname{TV}(w)=\sup\left\{
 \int_\R w(x)\phi'(x)\,\mathrm d\lambda_1(x):
 \phi\in C_c^1(\R),\ \|\phi\|_\infty\le1\right\}.
\]
Piszemy $w\in BV(\R)$, jeżeli $\operatorname{TV}(w)<\infty$.
Równoważnie pochodna dystrybucyjna $Dw$ jest skończoną
miarą ze znakiem, której wariacja całkowita jest równa
$\operatorname{TV}(w)$.
Normę miary o skończonej wariacji oznaczamy przez
$\|\mu\|_{\mathcal M}=|\mu|(\R)$.
Na $BV(\R)$ przyjmujemy normę
$\|w\|_{BV}=\|w\|_1+\operatorname{TV}(w)$.
\end{definicja}

\begin{uwaga}
$BV(\R)$ jest przestrzenią Banacha w tej normie.
Istotnie wariacja jest seminormą (wynika to z symetrii
zbioru testów i nierówności trójkąta).
Jeżeli $(w_n)$ jest ciągiem Cauchy'ego w normie $BV$,
to $w_n\to w$ w $L^1(\R)$. Dla każdego testu $\phi$
z definicji wariacji i każdego $n$ mamy
\[
 \left|\int_\R(w_n-w)\phi'\,\mathrm d\lambda_1\right|
 \le\liminf_{m\to\infty}\operatorname{TV}(w_n-w_m).
\]
Przechodząc do supremum po $\phi$, dostajemy to samo
oszacowanie dla $\operatorname{TV}(w_n-w)$.
Dla dużego $n$ prawa strona jest dowolnie mała.
W szczególności $w_n-w\in BV$, więc $w\in BV$,
a $\|w_n-w\|_{BV}\to0$.
\end{uwaga}

\begin{lemat}[Oszacowanie translacji w przestrzeni $BV$]
\label{lem:pde24-bv-translation}
Jeżeli $w\in BV(\R)$, to dla każdego $h\in\R$ zachodzi
\[
 \|w(\cdot+h)-w\|_1\le |h|\operatorname{TV}(w).
\]
\end{lemat}
\begin{proof}
Najpierw uzasadniamy równoważność z definicji przestrzeni
$BV$. Funkcjonał $\phi\mapsto-\int_\R w\phi'\,
\mathrm d\lambda_1$ na $C_c^1(\R)$ jest ograniczony
w normie supremum, z normą $\operatorname{TV}(w)$
(zbiór testów jest symetryczny). Rozszerza się
jednoznacznie na $C_0(\R)$ przez gęstość. Twierdzenie
Riesza o reprezentacji funkcjonałów na $C_0(\R)$ daje
skończoną miarę ze znakiem $Dw$ o tej samej normie
wariacji. Odwrotnie, jeżeli $Dw$ jest taką miarą,
to z definicji pochodnej dystrybucyjnej otrzymujemy
to samo ograniczenie, a gęstość testów daje równość norm.

Niech $\rho_\delta$ będzie nieujemnym gładkim jądrem
aproksymującym jedność i $w_\delta=\rho_\delta*w$.
Wtedy $w_\delta\to w$ w $L^1(\R)$ i
$w_\delta'=\rho_\delta*Dw$. Z twierdzenia Fubiniego
względem miary wariacji wynika
$\|w_\delta'\|_1\le |Dw|(\R)$.
Dla $h>0$ twierdzenie Newtona--Leibniza daje
\[
 \|w_\delta(\cdot+h)-w_\delta\|_1
 \le\int_0^h\int_\R|w_\delta'(x+s)|
 \,\mathrm d\lambda_1(x)\,\mathrm d\lambda_1(s)
 \le h\operatorname{TV}(w).
\]
Dla $h<0$ korzystamy z niezmienniczości normy względem
translacji, a dla $h=0$ nierówność jest trywialna.
Przechodzimy do granicy w $L^1$.
\end{proof}

\section{Skoki i warunek Rankine'a--Hugoniota}
\begin{twierdzenie}[Warunek skoku]
\label{thm:pde24-RH}
Niech $s\in C^1((0,T))$ i niech $u_-,u_+$ będą funkcjami
klasy $C^1$ odpowiednio po lewej i prawej stronie krzywej
$x=s(t)$, przy czym obie funkcje wraz z pierwszymi
pochodnymi przedłużają się ciągle do krzywej od swojej
strony. Oznaczamy ich ślady przez $u_L(t),u_R(t)$.
Załóżmy, że każda gałąź spełnia klasycznie
$u_t+\partial_xf(u)=0$. Funkcja sklejona z tych gałęzi
jest rozwiązaniem dystrybucyjnym w otoczeniu krzywej wtedy
i tylko wtedy, gdy dla każdego $t$ zachodzi
\begin{equation}\label{eq:pde24-RH}
 s'(t)\bigl(u_R(t)-u_L(t)\bigr)
 =f(u_R(t))-f(u_L(t)).
\end{equation}
\end{twierdzenie}
\begin{proof}
Niech $\varphi\in C_c^\infty((0,T)\times\R)$
ma nośnik w rozważanym otoczeniu. Z reguły Leibniza
dla całki po $x<s(t)$ i z równania klasycznego wynika
\[
 \int_0^T\!\int_{x<s(t)}(u_-\varphi_t+
 f(u_-)\varphi_x)\,\mathrm d\lambda_1(x)\,
 \mathrm d\lambda_1(t)
 =\int_0^T(f(u_L)-s'u_L)\varphi(t,s(t))\,
 \mathrm d\lambda_1(t).
\]
Istotnie człon czasowy daje po całkowaniu przez części
$-s'u_L\varphi$ na ruchomym brzegu, a człon przestrzenny
daje $f(u_L)\varphi$. Całki wewnętrzne znoszą się
na mocy równania klasycznego. Na obszarze $x>s(t)$
orientacja brzegu jest przeciwna, zatem otrzymujemy
analogicznie $(s'u_R-f(u_R))\varphi(t,s(t))$.
Po zsumowaniu i zmianie znaku, zgodnie z definicją
pochodnej dystrybucyjnej, dostajemy
\[
 \langle u_t+\partial_x f(u),\varphi\rangle
 =\int_0^T([f(u)]-s'[u])\varphi(t,s(t))\,
 \mathrm d\lambda_1(t),
 \qquad [h]=h_R-h_L.
\]
Prawą stronę zapisujemy także jako
$([f(u)]-s'[u])\delta_{x=s(t)}$, gdzie delta oznacza
miarę zadaną przez całkowanie względem czasu wzdłuż
wykresu, a nie względem jego długości łuku.
Współczynnik jest ciągły. Gdyby był niezerowy
w pewnym punkcie, miałby stały ścisły znak na pewnym
przedziale czasu; nieujemny test o nośniku
w małym otoczeniu odpowiedniego fragmentu wykresu
dałby niezerową całkę. Stąd znikanie dystrybucji
jest równoważne~\eqref{eq:pde24-RH}.
\end{proof}

Dla strumienia Burgersa $f(u)=u^2/2$, przy $u_L\ne u_R$,
prędkość skoku jest równa $(u_L+u_R)/2$. Sam warunek
Rankine'a--Hugoniota nie rozstrzyga jednak, czy taka fala
jest zgodna z fizycznym kierunkiem przepływu informacji.

\begin{kontrprzyklad}[Słabe rozwiązanie nie jest jednoznaczne]
\label{ex:pde24-expansion}
Dla $f(u)=u^2/2$ i danych Riemanna $u_0(x)=0$ dla $x<0$
oraz $u_0(x)=1$ dla $x>0$ funkcja
$u(t,x)=0$ dla $x<t/2$ i $u(t,x)=1$ dla $x>t/2$
spełnia warunek Rankine'a--Hugoniota, więc jest rozwiązaniem
słabym. Innym rozwiązaniem słabym tych samych danych jest
ciągły wachlarz $u(t,x)=0$ dla $x\le0$,
$u(t,x)=x/t$ dla $0<x<t$ oraz $u(t,x)=1$ dla $x\ge t$.
Na wnętrzu wachlarza $u_t=-x/t^2$ oraz
$\partial_x(u^2/2)=x/t^2$; na pozostałych
obszarach funkcja jest stała. Na prostych $x=0$ i
$x=t$ ślady obu gałęzi są równe, więc warunek
~\eqref{eq:pde24-RH} jest spełniony. Dla obu
rozwiązań różnica od $u_0$ jest ograniczona przez
$1$ i znika poza $[0,t]$, zatem norma tej różnicy
w $L^1_{\mathrm{loc}}$ dąży do zera przy $t\downarrow0$.
Całkowanie przez części od dodatniej chwili i przejście
z tą chwilą do zera daje pełny wzór
~\eqref{eq:pde24-weak}. Dodatkowa zasada selekcji
jest konieczna.
\end{kontrprzyklad}

\section{Entropie i poprawny kierunek fal uderzeniowych}
\begin{definicja}[Para entropijna]
Niech $f\in C^1(\R)$. Parą entropijną nazywamy funkcje
$\eta\in C^2(\R)$ oraz $q\in C^1(\R)$ takie, że $\eta$ jest
wypukła i $q'(z)=\eta'(z)f'(z)$. Rozwiązanie spełnia
nierówność entropijną, jeśli
$\partial_t\eta(u)+\partial_x q(u)\le0$ w sensie dystrybucji.
\end{definicja}

\begin{lemat}[Produkcja entropii przez lepkość]
\label{lem:pde24-viscentropy}
Jeżeli $u^\eps$ jest rozwiązaniem klasy $C^{1,2}$
(jedna ciągła pochodna czasowa i dwie przestrzenne)
równania $u_t^\eps+\partial_xf(u^\eps)=\eps u_{xx}^\eps$,
$\eps>0$, to dla każdej pary entropijnej zachodzi tożsamość
\begin{equation}\label{eq:pde24-entropyvisc}
 \partial_t\eta(u^\eps)+\partial_xq(u^\eps)
 =\eps\partial_{xx}\eta(u^\eps)
 -\eps\eta''(u^\eps)|u_x^\eps|^2
 \le\eps\partial_{xx}\eta(u^\eps).
\end{equation}
\end{lemat}
\begin{proof}
Mnożymy równanie przez $\eta'(u^\eps)$ i używamy reguły
łańcuchowej. Z równości
$\eta'(u)u_{xx}=\partial_{xx}\eta(u)-\eta''(u)u_x^2$
wynika wzór; znak ostatniego członu wynika z wypukłości.
\end{proof}

\begin{definicja}[Rozwiązanie entropijne Krużkowa]
\label{def:pde24-kru}
Niech $u_0\in L^1(\R)\cap L^\infty(\R)$.
Funkcja $u\in C([0,T];L^1(\R))\cap L^\infty((0,T)\times\R)$
jest \emph{rozwiązaniem entropijnym}~\eqref{eq:pde24-cl}, jeżeli
$u(0)=u_0$ w $L^1$ i dla każdego $k\in\R$ zachodzi
\begin{equation}\label{eq:pde24-kruzineq}
 \partial_t|u-k|+\partial_x\left(
   \operatorname{sgn}(u-k)(f(u)-f(k))\right)\le0
 \quad\text{w }\mathcal D'((0,T)\times\R),
\end{equation}
gdzie $\operatorname{sgn}(0)=0$.
Oznaczając $q_k(z)=\operatorname{sgn}(z-k)(f(z)-f(k))$,
zapisujemy tę nierówność równoważnie jako
\[
 \int_0^T\!\int_\R
 \bigl(|u-k|\varphi_t+q_k(u)\varphi_x\bigr)
 \,\mathrm d\lambda_1(x)\,\mathrm d\lambda_1(t)\ge0
\]
dla każdej nieujemnej funkcji
$\varphi\in C_c^\infty((0,T)\times\R)$.
\end{definicja}

\begin{uwaga}
Z warunków dla wszystkich $k$
wynika również słabe równanie: jeśli $k$ jest większe od
istotnego supremum $u$, w~\eqref{eq:pde24-kruzineq}
uzyskujemy jedną nierówność liniową, a jeśli jest mniejsze
od istotnego infimum, nierówność przeciwną. Obie nierówności dają równość dystrybucyjną.
Aby odzyskać również~\eqref{eq:pde24-weak},
testujemy równanie przez $\varphi(t,x)\theta_\delta(t)$,
gdzie $\theta_\delta$ rośnie gładko od $0$ do $1$
w $(\delta,2\delta)$ i jest stała poza tym przedziałem.
Wyraz zawierający $\theta_\delta'$ dąży do
$\int_\R u_0(x)\varphi(0,x)\,\mathrm d\lambda_1(x)$:
pochodna ma masę $1$, a całka z $u(t,x)\varphi(t,x)$
jest ciągła w chwili $0$. Pozostałe wyrazy przechodzą
do granicy przez zbieżność zmajoryzowaną na nośniku
testu. Tak samo postępujemy przy lokalnym śladzie
w $L^1_{\mathrm{loc}}$; przy lokalnej ograniczoności
stałe $k$ dobieramy stosownie do nośnika testu.
\end{uwaga}

\begin{lemat}[Równoważność warunków entropijnych]
\label{lem:pde24-entropy-equivalence}
Niech $u$ będzie lokalnie ograniczonym rozwiązaniem
dystrybucyjnym równania zachowania na $(0,T)\times\R$.
Nierówności Krużkowa~\eqref{eq:pde24-kruzineq} dla
wszystkich $k\in\R$ są równoważne nierównościom
$\partial_t\eta(u)+\partial_xq(u)\le0$ dla wszystkich
gładkich par entropijnych w sensie podanej wyżej definicji.
\end{lemat}
\begin{proof}
Ustalmy nieujemny test $\varphi$ i przedział $[A,B]$
zawierający istotny zakres wartości $u$ na jego nośniku.
Jeżeli zachodzą nierówności Krużkowa, to dla wypukłej
$\eta\in C^2(\R)$ na $[A,B]$ mamy
\[
 \eta(z)=a+bz+\frac12\int_A^B |z-k|\eta''(k)\,
                         \mathrm d\lambda_1(k),
\]
gdzie $a,b$ są stałymi. Istotnie pochodna rzędu
drugiego prawej strony na $(A,B)$ jest równa
$\eta''(z)$, więc różnica jest afiniczna.
Dla $q'=\eta'f'$ analogicznie
\[
 q(z)=c+bf(z)+\frac12\int_A^B q_k(z)\eta''(k)\,
                         \mathrm d\lambda_1(k).
\]
Druga tożsamość wynika z różniczkowania względem
$z$: dla $k\ne z$ pochodna $q_k(z)$ wynosi
$\operatorname{sgn}(z-k)f'(z)$, a ilorazy różnicowe
są jednostajnie ograniczone przez supremum $|f'|$
na przedziale. Twierdzenie o zbieżności zmajoryzowanej
pozwala różniczkować pod całką; stała $c$ uwzględnia
dowolną normalizację strumienia.
Części afiniczne dają zerowy wkład do nierówności,
ponieważ $u_t+\partial_x f(u)=0$ dystrybucyjnie.
Pozostały wkład jest całką z nierówności Krużkowa
z nieujemną wagą $\eta''/2$. Twierdzenie Fubiniego
stosujemy na zwartym nośniku testu i ograniczonym
przedziale parametrów $k$. Otrzymujemy żądaną
nierówność dla $\eta,q$.

Odwrotnie, dla ustalonego $k$ bierzemy
$\eta_\delta(z)=\sqrt{(z-k)^2+\delta^2}$ oraz
$q_\delta(z)=\int_k^z\eta_\delta'(r)f'(r)\,
\mathrm d\lambda_1(r)$. Są to gładkie pary entropijne
w wymaganym sensie. Entropie zbiegają jednostajnie
do $|z-k|$, a strumienie do $q_k$ jednostajnie na
$[A,B]$: błąd strumienia nie przekracza całki
z $|\eta_\delta'(r)-\operatorname{sgn}(r-k)|\,|f'(r)|$
po przedziale zawierającym $[A,B]$ i $k$.
Ta całka dąży do zera przez zbieżność zmajoryzowaną.
Przechodząc do granicy w całkach z pochodnymi testu,
otrzymujemy nierówność Krużkowa.
\end{proof}

\begin{twierdzenie}[Warunek entropijny na skoku]
\label{thm:pde24-entropyshock}
Niech $u_L\ne u_R$, $f\in C^2(\R)$ i niech skok między
stałymi stanami $u_L,u_R$ porusza się z prędkością
$\sigma=(f(u_R)-f(u_L))/(u_R-u_L)$.
Jest on entropijny w sensie lokalnych nierówności
Krużkowa wtedy i tylko wtedy, gdy dla każdego $k$ pomiędzy
$u_L$ i $u_R$ zachodzi
\begin{equation}\label{eq:pde24-chord}
 \begin{cases}
 f(k)\le \ell(k),&u_L>u_R,\\
 f(k)\ge \ell(k),&u_L<u_R,
 \end{cases}
 \qquad
 \ell(k)=f(u_L)+\sigma(k-u_L).
\end{equation}
W szczególności gdy $f$ jest ściśle wypukła, skok z
$u_L>u_R$ jest entropijny, natomiast skok z $u_L<u_R$
nie jest entropijny.
\end{twierdzenie}
\begin{proof}
Przyjmujemy $[h]=h_R-h_L$. Z rachunku dystrybucyjnego
nierówność dla entropii $\eta$ jest równoważna
$[q]-\sigma[\eta]\le0$. Jeżeli $k$ leży poza odcinkiem
wyznaczonym przez stany, wyrażenie to znika, ponieważ
$[q]=\pm[f]$ i $[\eta]=\pm[u]$.
Gdy $u_L>u_R$ i $u_R<k<u_L$, mamy
$[q_k]=2f(k)-f(u_R)-f(u_L)$ oraz
$[|u-k|]=2k-u_R-u_L$. Stąd
\[
 \begin{aligned}
 [q_k]-\sigma[|u-k|]
 &=2f(k)-f(u_R)-f(u_L)\\
 &\quad-\sigma(2k-u_R-u_L)\\
 &=2(f(k)-\ell(k)),
 \end{aligned}
\]
gdzie w ostatniej równości korzystamy z
$f(u_R)=f(u_L)+\sigma(u_R-u_L)$.
Przy $u_L<u_R$ i $u_L<k<u_R$ oba wyrażenia
mają przeciwne znaki, więc wynik to
$2(\ell(k)-f(k))$. Dla $k$ równego jednemu ze
stanów współczynnik jest zerowy przez ciągłość.
Te warunki są konieczne i wystarczające. Ścisła wypukłość
umieszcza wykres $f$ ściśle poniżej cięciwy łączącej
odpowiednie punkty, co kończy dowód.
\end{proof}

\section{Zagadnienie Riemanna dla ściśle wypukłego strumienia}
\begin{twierdzenie}[Fala uderzeniowa i rozrzedzenia]
\label{thm:pde24-riemann}
Niech $f\in C^2(\R)$ i $f''(z)>0$ na domkniętym
przedziale o końcach $u_L,u_R$, jeśli $u_L\ne u_R$. Rozważmy dane $u_0=u_L$ dla $x<0$ oraz
$u_0=u_R$ dla $x>0$ (nie muszą należeć do $L^1$).
Dla $t>0$ istnieje samopodobne lokalnie ograniczone
rozwiązanie entropijne $u(t,x)=U(x/t)$:
\begin{enumerate}
\item Jeśli $u_L>u_R$, jest to skok przemieszczający się
z prędkością $\sigma=(f(u_L)-f(u_R))/(u_L-u_R)$.
\item Jeśli $u_L<u_R$, to
\[
 U(\xi)=\begin{cases}
 u_L,&\xi\le f'(u_L),\\
 (f')^{-1}(\xi),&f'(u_L)<\xi<f'(u_R),\\
 u_R,&\xi\ge f'(u_R).
 \end{cases}
\]
\end{enumerate}
Jeżeli $u_L=u_R$, rozwiązaniem jest funkcja stała $u(t,x)=u_L$.
Rozwiązania te są jednoznaczne w klasie
$C([0,T];L^1_{\mathrm{loc}}(\R))\cap L^\infty((0,T)\times\R)$
rozwiązań lokalnych nierówności entropijnych z zadanym
śladem początkowym. Dowód wykorzystuje skończoną prędkość
propagacji (Wniosek~\ref{cor:pde24-finite}).
Nie zakładamy całkowalności na całej prostej, lecz
wymagamy globalnego ograniczenia wartości rozwiązania.
\end{twierdzenie}
\begin{proof}
W pierwszym przypadku warunek Rankine'a--Hugoniota
sprawia, że równanie jest spełnione słabo, a
Twierdzenie~\ref{thm:pde24-entropyshock} zapewnia
nierówności entropijne. W drugim przypadku $U$ jest
ciągła, a na wachlarzu spełnia $f'(U(\xi))=\xi$.
Obliczamy $u_t=-\xi U'(\xi)/t$ oraz
$\partial_x f(u)=f'(U)U'(\xi)/t$; składniki się znoszą.
Poza wachlarzem rozwiązanie jest stałe. Ponieważ funkcja
jest ciągła, w przejściu przez brzegi wachlarza nie
powstają wyrazy typu delta. Dla każdej gładkiej pary entropijnej
na wnętrzu wachlarza zachodzi równość
$\partial_t\eta(u)+\partial_xq(u)=0$.
Po uzasadnieniu braku składników osobliwych poniżej
stosujemy Lemat~\ref{lem:pde24-entropy-equivalence}.
W obu przypadkach $u(t,\cdot)$ jest globalnie ograniczona
przez $\max(|u_L|,|u_R|)$. Różnica od danych początkowych
może wystąpić tylko w $[-Ct,Ct]$, gdzie
$C=\max_{z\in[\min(u_L,u_R),\max(u_L,u_R)]}|f'(z)|$
(zaś $|\sigma|\le C$ przez twierdzenie o wartości średniej).
Jej norma $L^1$ na dowolnym ograniczonym przedziale
dąży więc do zera. Dla $t_n\to t>0$ profile zbiegają
punktowo poza ewentualnym punktem skoku, więc zbieżność
zmajoryzowana daje ciągłość w $L^1_{\mathrm{loc}}$.
W przypadku wachlarza założenie o dodatniości $f''$
na domkniętym przedziale gwarantuje ograniczoność
pochodnej $(f')^{-1}$; na każdym zwartym przedziale
dodatnich czasów profil jest kawałkami klasy $C^1$
i ciągły na granicach kawałków. Całkowanie przez części
jak w dowodzie warunku skoku pokazuje, że nie ma
składników osobliwych także dla $\eta(u),q(u)$.
Jednoznaczność wynika z Wniosku~\ref{cor:pde24-finite};
dla równych stanów stosujemy rozwiązanie stałe.
\end{proof}

\begin{center}
\begin{tikzpicture}[x=1.05cm,y=1.55cm,>=Latex]
 \begin{scope}
  \draw[->] (-2.15,0)--(2.15,0) node[right] {$x$};
  \draw[->] (0,-.08)--(0,1.9) node[above] {$t$};
  \foreach \a in {-1.8,-1.2,-.6}{
   \pgfmathsetmacro{\ttt}{min(1.65,-2*\a)}
   \draw[blue!70!black] (\a,0)--({\a+\ttt},\ttt);
  }
  \foreach \a in {.3,.6,.9,1.2,1.5,1.8}{
   \pgfmathsetmacro{\ttt}{min(1.65,2*\a)}
   \draw[teal!70!black] (\a,0)--(\a,\ttt);
  }
  \draw[red!75!black,very thick] (0,0)--(.825,1.65)
     node[above right,font=\small] {$x=t/2$};
  \node[font=\small] at (-1.25,1.73) {$u_L=1$};
  \node[font=\small] at (1.55,1.73) {$u_R=0$};
  \node[font=\small] at (0,-.35) {fala uderzeniowa};
 \end{scope}
 \begin{scope}[xshift=5.1cm]
  \draw[->] (-2.15,0)--(2.15,0) node[right] {$x$};
  \draw[->] (0,-.08)--(0,1.9) node[above] {$t$};
  \foreach \a in {-1.8,-1.2,-.6}
   \draw[blue!70!black] (\a,0)--(\a,1.65);
  \foreach \a in {.3,.6,.9,1.2,1.5,1.8}{
   \pgfmathsetmacro{\ttt}{min(1.65,2.1-\a)}
   \draw[teal!70!black] (\a,0)--({\a+\ttt},\ttt);
  }
  \foreach \v in {.2,.4,.6,.8}
   \draw[orange!80!black] (0,0)--({1.65*\v},1.65);
  \draw[orange!80!black,very thick] (0,0)--(1.65,1.65)
      node[above right,font=\small] {$x=t$};
  \node[font=\small] at (-1.25,1.73) {$u_L=0$};
  \node[font=\small,anchor=west] at (.7,-.14) {$u_R=1$};
  \node[font=\small] at (0,-.35) {fala rozrzedzenia};
 \end{scope}
\end{tikzpicture}

\small Rysunek: charakterystyki równania Burgersa dla
danych $(u_L,u_R)=(1,0)$ (po lewej) oraz $(0,1)$
(po prawej). W pierwszym przypadku charakterystyki
dochodzą do skoku o prędkości $1/2$; w drugim
wachlarz promieni $x=vt$, $0<v<1$, wypełnia
obszar między charakterystykami $x=0$ i $x=t$.
\end{center}

\section{Nierówność Kato dla dwóch rozwiązań entropijnych}
\begin{twierdzenie}[Nierówność Kato--Krużkowa]
\label{thm:pde24-kato}
Niech $f\in C^1(\R)$ oraz niech
$u,v\in C([0,T];L^1(\R))\cap L^\infty((0,T)\times\R)$
będą rozwiązaniami entropijnymi z danymi $u_0,v_0$.
Wówczas w sensie dystrybucji na $(0,T)\times\R$
\begin{equation}\label{eq:pde24-kato}
 \partial_t|u-v|+\partial_x\bigl[
 \operatorname{sgn}(u-v)(f(u)-f(v))\bigr]\le0.
\end{equation}
\end{twierdzenie}
\begin{proof}
\textbf{Krok 1: podwojenie zmiennych.}
Dla ustalonej pary $(s,y)$ w nierówności entropijnej
dla $u(t,x)$ wybieramy $k=v(s,y)$, a dla ustalonej pary
$(t,x)$ w nierówności dla $v(s,y)$ wybieramy $k=u(t,x)$.
Są to dopuszczalne stałe, ponieważ w każdej nierówności
całkujemy tylko po zmiennych jej rozwiązania. Po
scalaniu nierówności twierdzeniem Fubiniego otrzymujemy
sumę dwóch nierówności na $(t,x,s,y)$.

\textbf{Krok 2: funkcja testowa.}
Niech $\rho\in C_c^\infty(\R)$ jest nieujemna, parzysta
i ma całkę równą $1$. Dla $\delta>0$ kładziemy
$\rho_\delta(z)=\delta^{-1}\rho(z/\delta)$.
Bierzemy $\psi\in C_c^\infty((0,T)\times\R)$, $\psi\ge0$,
i testujemy obie nierówności funkcją
\[
 \Phi_\delta(t,x,s,y)=
 \rho_\delta(t-s)\rho_\delta(x-y)\,
 \psi\!\left(\frac{t+s}{2},\frac{x+y}{2}\right).
\]
Dla małego $\delta$ nośnik leży w zwartym podzbiorze
$((0,T)\times\R)^2$: średnie zmiennych należą do
nośnika $\psi$, a ich różnice do nośników jąder.
Zatem dla każdej ustalonej pary zmiennych drugi
przekrój jest dopuszczalnym testem. Wybieramy mierzalne
reprezentanty $u,v$; wszystkie funkcje podcałkowe są
mierzalne, ograniczone na tym zwartym nośniku
i całkowalne, więc stosowanie twierdzenia Fubiniego
jest uzasadnione. Warunek entropijny obowiązuje dla
każdej stałej, nie tylko dla prawie każdej.
Suma pochodnych spełnia dokładnie
\[
 (\partial_t+\partial_s)\Phi_\delta
 =\rho_\delta(t-s)\rho_\delta(x-y)\psi_t,
 \qquad
 (\partial_x+\partial_y)\Phi_\delta
 =\rho_\delta(t-s)\rho_\delta(x-y)\psi_x,
\]
gdzie pochodne $\psi$ obliczamy w zmiennych średnich.
Po dodaniu i zastosowaniu twierdzenia Fubiniego dostajemy
\[
 \int_{((0,T)\times\R)^2}
 \left(
 |u(t,x)-v(s,y)|(\partial_t+\partial_s)\Phi_\delta
 +Q(u(t,x),v(s,y))(\partial_x+\partial_y)\Phi_\delta
 \right)
 \,\mathrm d\lambda_4(t,x,s,y)\ge0,
\]
gdzie $Q(a,b)=\operatorname{sgn}(a-b)(f(a)-f(b))$.
W obu nierównościach występują te same wartości
$|u(t,x)-v(s,y)|$ i ten sam strumień
$\operatorname{sgn}(u(t,x)-v(s,y))
(f(u(t,x))-f(v(s,y)))$; w nierówności dla $v$
oba znaki zmieniają się jednocześnie.

\textbf{Krok 3: granica diagonalna.}
Podstawiamy $h=t-s$, $k=x-y$; wartość bezwzględna
jakobianu tego podstawienia wynosi $1$. Wówczas
$v(s,y)=v(t-h,x-k)$, a pochodne $\psi$ obliczamy
w $(t-h/2,x-k/2)$. Na ustalonym zwartym powiększeniu
nośnika $\psi$ ciągłość translacji w $L^1_{\mathrm{loc}}$
daje
\[
 \int |v(t-h,x-k)-v(t,x)|\,\mathrm d\lambda_2(t,x)
 \longrightarrow0
 \qquad ((h,k)\to(0,0)),
\]
gdzie całka jest po tym powiększeniu. Różnicę wartości
$|u-v|$ ograniczamy przez $|v(t-h,x-k)-v(t,x)|$.
Jeżeli $M$ ogranicza $|f'|$ na wspólnym przedziale wartości
$u,v$, to dla ustalonego $a$ funkcja $b\mapsto Q(a,b)$
jest $M$-lipschitzowska: na każdej stronie punktu $a$
wynika to z twierdzenia o wartości średniej, a dla
argumentów po przeciwnych stronach stosujemy nierówność
trójkąta i $Q(a,a)=0$. Tak samo postępujemy w pierwszym
argumencie. Błąd w strumieniu ma więc analogiczne
oszacowanie. Całkowanie względem nieujemnego jądra
$\rho_\delta(h)\rho_\delta(k)$ o masie $1$ nie zwiększa
supremum tych błędów po jego nośniku. Pochodne $\psi$
w przesuniętych argumentach zbiegają jednostajnie
do pochodnych w $(t,x)$, a pozostałe czynniki są
ograniczone na wspólnym zwartym nośniku. Ostatecznie
\[
 \int_0^T\!\int_\R
 \bigl(|u-v|\psi_t+
 \operatorname{sgn}(u-v)(f(u)-f(v))\psi_x\bigr)
 \,\mathrm d\lambda_1(x)\,\mathrm d\lambda_1(t)\ge0.
\]
Jest to~\eqref{eq:pde24-kato} w sensie dystrybucyjnym.
\end{proof}

\begin{twierdzenie}[Kontrakcja w $L^1$ i jednoznaczność]
\label{thm:pde24-L1}
Przy założeniach Twierdzenia~\ref{thm:pde24-kato} zachodzi
\begin{equation}\label{eq:pde24-L1}
 \|u(t)-v(t)\|_{L^1(\R)}
 \le\|u_0-v_0\|_{L^1(\R)},\qquad 0\le t\le T.
\end{equation}
W szczególności dane początkowe wyznaczają co najwyżej
jedno rozwiązanie entropijne w powyższej klasie.
\end{twierdzenie}
\begin{proof}
Oznaczmy $w=|u-v|\in C([0,T];L^1)$ i
$q=\operatorname{sgn}(u-v)(f(u)-f(v))$. Jeśli
$M=\max_{|z|\le B}|f'(z)|$, gdzie $B$ ogranicza
$|u|,|v|$, to $|q|\le Mw$. W nierówności Kato
testujemy nieujemną funkcją aproksymującą
$\chi_{[s,t]}(\tau)\zeta_R(x)$, gdzie $\zeta_R$ jest
równa $1$ na $[-R,R]$, równa $0$ poza $[-2R,2R]$
i spełnia $|\zeta_R'|\le C/R$. Przechodząc do granicy
w aproksymacji czasowej, otrzymujemy
\[
 \int_\R w(t,x)\zeta_R(x)\,\mathrm d\lambda_1(x)
 \le \int_\R w(s,x)\zeta_R(x)\,\mathrm d\lambda_1(x)
 +\frac{CM}{R}\int_s^t\!\int_\R w(\tau,x)
 \,\mathrm d\lambda_1(x)\,\mathrm d\tau.
\]
Wybieramy $0\le\zeta_R\le1$. Dla $R\to\infty$
całki po obu stronach przechodzą do norm $L^1$
na mocy twierdzenia o zbieżności zmajoryzowanej.
Ostatni człon znika, ponieważ $w$ jest ciągła
w $L^1$ na zwartym przedziale czasowym.
Otrzymujemy kontrakcję dla $0<s<t<T$,
a przypadki $s=0$ i $t=T$ wynikają z ciągłości
w $L^1$.
\end{proof}

\begin{lemat}[Lokalna postać nierówności Kato]
\label{lem:pde24-local-kato}
Niech $f\in C^1(\R)$, a funkcje $u,v$ niech należą do
\[
 C([0,T];L^1_{\mathrm{loc}}(\R))\cap
 L^\infty_{\mathrm{loc}}([0,T]\times\R).
\]
Załóżmy, że spełniają lokalnie wszystkie nierówności entropijne
Krużkowa i mają lokalne ślady początkowe $u_0,v_0$.
Wówczas nierówność~\eqref{eq:pde24-kato} zachodzi
w $\mathcal D'((0,T)\times\R)$. Do wyprowadzenia
lokalnej kontrakcji na zwężających się przedziałach
wystarcza ograniczenie $|f'|$ na zakresie wartości
$u,v$ w rozważanym stożku.
\end{lemat}
\begin{proof}
Każda funkcja testowa $\psi\in C_c^\infty((0,T)\times\R)$
ma nośnik wewnątrz pewnego zwartego walca.
W dowodzie Twierdzenia~\ref{thm:pde24-kato}
całkujemy cztery zmienne tylko w niewielkim powiększeniu
tego walca, bo funkcje
$\rho_\delta(t-s)\rho_\delta(x-y)$ mają małe nośniki.
Zatem potrzebujemy jedynie ograniczeń $L^\infty$
i ciągłości translacji w $L^1$ na tym walcu,
a nie całkowalności $u,v$ na całej prostej.
Przejście diagonalne daje nierówność Kato lokalnie.
W testowaniu przez przybliżenie wskaźnika stożka
stosujemy analogiczne odcięcie przestrzenne; używane
granice czasu wynikają z założonej ciągłości w
$L^1_{\mathrm{loc}}$. Warunek $|f'|\le M$ na
wspólnym zakresie wartości zapewnia $|q|\le M|u-v|$.
\end{proof}

\begin{wniosek}[Skończona prędkość propagacji]
\label{cor:pde24-finite}
Niech $u,v$ spełniają założenia
Lematu~\ref{lem:pde24-local-kato}. Ustalmy $a<b$,
$t\in[0,T]$ oraz $M\ge0$. Załóżmy, że $|f'|\le M$
na przedziale zawierającym istotne zakresy wartości
$u,v$ w zbiorze
\[
 \{(\tau,x):0\le\tau\le t,\ 
 a-M(t-\tau)\le x\le b+M(t-\tau)\}
\]
i jego pewnym otoczeniu względnym w $[0,T]\times\R$.
Jeżeli $u,v$ są globalnie ograniczone przez $B$,
to $M=\max_{|z|\le B}|f'(z)|$ działa dla wszystkich
$a,b,t$. Przy podanych założeniach zachodzi
\begin{equation}\label{eq:pde24-cone}
 \int_a^b|u(t,x)-v(t,x)|\,\mathrm d\lambda_1(x)
 \le\int_{a-Mt}^{b+Mt}|u_0(x)-v_0(x)|\,\mathrm d\lambda_1(x).
\end{equation}
\end{wniosek}
\begin{proof}
Na mocy Lematu~\ref{lem:pde24-local-kato}
możemy zastosować nierówność Kato lokalnie. Używamy
funkcji testowych przybliżających wskaźnik zwężającego się przedziału
$[a-M(t-\tau),b+M(t-\tau)]$ dla $0\le\tau\le t$.
Jego lewy koniec $L(\tau)=a-M(t-\tau)$ porusza się
z prędkością $M$, a prawy $R(\tau)=b+M(t-\tau)$
z prędkością $-M$. Niech $\rho_\delta$ będzie nieujemnym gładkim jądrem
o całce $1$ i nośniku w $[-\delta,\delta]$. Zamiast
podstawiać niegładką funkcję charakterystyczną lub
wymagać istnienia śladów $w=|u-v|$ i $q$ na brzegach
stożka, definiujemy
\[
 \zeta_\delta(\tau,x)=
 \int_{L(\tau)}^{R(\tau)}\rho_\delta(x-z)\,
 \mathrm d\lambda_1(z).
\]
Z reguły Leibniza oraz całkowania po $z$ wynika
\[
 \partial_\tau\zeta_\delta
 =-M\bigl(\rho_\delta(x-L(\tau))+
             \rho_\delta(x-R(\tau))\bigr),
 \qquad
 \partial_x\zeta_\delta
 =\rho_\delta(x-L(\tau))-\rho_\delta(x-R(\tau)).
\]
Stąd prawie wszędzie
\[
 w\partial_\tau\zeta_\delta+q\partial_x\zeta_\delta
 =(q-Mw)\rho_\delta(x-L(\tau))
   +(-q-Mw)\rho_\delta(x-R(\tau))\le0.
\]
Mnożymy $\zeta_\delta$ przez nieujemne gładkie
odcięcie czasowe przybliżające $\chi_{[s,t]}$,
gdzie $0<s<t<T$, i stosujemy nierówność Kato.
Po przejściu do granicy w odcięciu czasowym otrzymujemy
\[
 \int_\R w(t,x)\zeta_\delta(t,x)\,\mathrm d\lambda_1(x)
 \le
 \int_\R w(s,x)\zeta_\delta(s,x)\,\mathrm d\lambda_1(x).
\]
Ciągłość $w$ w $L^1_{\mathrm{loc}}$ pozwala
przejść z $s$ do $0$, a z $t$ do $T$, gdy jest to
potrzebne. Ponieważ $0\le\zeta_\delta\le1$,
nośniki pozostają w jednym ograniczonym przedziale,
a $\zeta_\delta$ zbiega prawie wszędzie do funkcji
charakterystycznej odpowiedniego przedziału,
twierdzenie o zbieżności zmajoryzowanej przy
$\delta\downarrow0$ daje~\eqref{eq:pde24-cone}.
Dla $t=0$ nierówność jest równością.
\end{proof}

\section{Istnienie przez zanikającą lepkość}
\begin{twierdzenie}[Istnienie dla danych $BV$]
\label{thm:pde24-existence}
Niech $f\in C^2(\R)$ oraz
$u_0\in BV(\R)\cap L^1(\R)\cap L^\infty(\R)$.
Wtedy dla każdego $T>0$ istnieje dokładnie jedno
rozwiązanie entropijne $u\in C([0,T];L^1(\R))\cap L^\infty((0,T)\times\R)$
z danymi $u_0$. Dla każdego $t\in[0,T]$ spełnia ono
\begin{equation}\label{eq:pde24-BVbound}
 \|u(t)\|_\infty\le\|u_0\|_\infty,
 \qquad \operatorname{TV}(u(t))\le\operatorname{TV}(u_0),
 \qquad \|u(t)\|_1\le\|u_0\|_1.
\end{equation}
\end{twierdzenie}
\begin{proof}
\textbf{Krok 1: aproksymacja lepkościowa.}
Zastępujemy $f$ przez $f-f(0)$, co nie zmienia równania,
i oznaczamy nadal tę funkcję przez $f$.
Niech $u_0^\eps=G_{\eps^2}*u_0$, gdzie $G_{\eps^2}$ jest
nieujemnym gładkim jądrem o masie $1$ i szerokości
$O(\eps)$. Rozwiązujemy
\begin{equation}\label{eq:pde24-visc}
 u_t^\eps+\partial_x f(u^\eps)=\eps u_{xx}^\eps,
 \qquad u^\eps(0)=u_0^\eps.
\end{equation}
Dla ustalonego $\eps>0$ rozwiązanie otrzymujemy
lokalnie z wzoru Duhamela
\[
 u^\eps(t)=G_{\eps t}*u_0^\eps
 -\int_0^t\partial_xG_{\eps(t-s)}*f(u^\eps(s))\,\mathrm ds,
\]
gdzie $G_t(x)=(4\pi t)^{-1/2}e^{-x^2/(4t)}$.
Uzasadnimy istnienie i całkowalność jeszcze przed
użyciem kontrakcji lepkościowej. Ustalmy $\eps>0$ i
$B=\|u_0\|_\infty$. Zastępujemy tymczasowo $f$ przez
funkcję $\widetilde f\in C^2(\R)$ o ograniczonych
pierwszej i drugiej pochodnej, równą $f$ na $[-B,B]$,
na przykład mnożąc $f$ przez gładkie odcięcie równe
$1$ na otoczeniu tego przedziału. Niech
$L=\|\widetilde f'\|_\infty$.
W przestrzeni Banacha
$X=L^1(\R)\cap BUC(\R)$, z normą
$\|w\|_X=\|w\|_1+\|w\|_\infty$, przez $BUC$ rozumiemy
ograniczone funkcje jednostajnie ciągłe. Zupełność
$X$ wynika z zupełności obu przestrzeni: granice ciągu
Cauchy'ego w $L^1$ oraz w normie supremum są równe
prawie wszędzie, co widać po wyborze podciągu
zbieżnego prawie wszędzie w $L^1$. Jądro Gaussa
jest silnie ciągłą półgrupą kontrakcji na $X$:
w normie $L^1$ wynika to z ciągłości translacji,
a w normie supremum z jednostajnej ciągłości funkcji.
Ponieważ $\widetilde f(0)=0$, superpozycja
$w\mapsto\widetilde f(w)$ jest $L$-lipschitzowska na $X$.
Dla prawej strony wzoru Duhamela, oznaczonej przez
$\mathcal P$, nierówność Younga i oszacowanie
$\|\partial_xG_{\eps t}\|_1=C(\eps t)^{-1/2}$ dają
\[
 \|\mathcal P a-\mathcal P b\|_{C([0,S];X)}
 \le 2CL\sqrt{S/\eps}\,
     \|a-b\|_{C([0,S];X)}.
\]
Całka definiuje funkcję ciągłą w czasie: osobno
rozpatrujemy część oddaloną od $s=t$, gdzie jądro
zależy ciągle od czasu w $L^1$, oraz część bliską
$s=t$, ograniczoną przez całkę z $(t-s)^{-1/2}$.
Dla małego $S$ stała kontrakcji jest mniejsza od $1$.
Twierdzenie Banacha o punkcie stałym i powtarzanie
konstrukcji na przedziałach tej samej długości daje
rozwiązanie w $C([0,T];X)$. Ten sam argument
w samej przestrzeni $BUC(\R)$ daje jednoznaczność
rozwiązania całkowego także w $C([0,T];BUC(\R))$:
stała kontrakcji w normie supremum jest taka sama.
Będziemy korzystać z tej jednoznaczności przy
uzasadnianiu dodatkowej regularności.

Dla gładkich danych Gaussa rozwiązanie ma także
regularność klasyczną. Można to uzasadnić lokalnie
tym samym argumentem w przestrzeni funkcji klasy
$C^1$, dla których funkcja i jej pochodna są
ograniczone oraz jednostajnie ciągłe, z normą
$\|w\|_\infty+\|w_x\|_\infty$. Jest to przestrzeń
Banacha: jednostajne granice funkcji i pochodnych
są związane twierdzeniem Newtona--Leibniza.
Stosujemy odwzorowanie Picarda na domkniętej kuli
wokół ewolucji danych przez półgrupę Gaussa;
dla małego czasu pozostaje ono w tej kuli i jest
kontrakcją, gdyż całka z normy pochodnej jądra
jest rzędu $\sqrt{S/\eps}$. Dokładniej, po zróżniczkowaniu wzoru Duhamela
przenosimy jedną pochodną na strumień, czyli używamy
$\partial_x\widetilde f(w)=\widetilde f'(w)w_x$.
Dla funkcji $a,b$ z ograniczonej kuli
$\|a_x\|_\infty,\|b_x\|_\infty\le K$ mamy
\[
 \begin{aligned}
 \|(\widetilde f(a)-\widetilde f(b))_x\|_\infty
 &\le L\|a_x-b_x\|_\infty\\
 &\quad+K\|\widetilde f''\|_\infty\|a-b\|_\infty .
 \end{aligned}
\]
Wraz z oszacowaniem w normie supremum daje to
lokalną lipschitzowskość superpozycji w tej normie.
Otrzymany punkt stały jest równy wcześniej
skonstruowanemu rozwiązaniu dzięki jednoznaczności
w $C([0,S];BUC(\R))$. Dla $z(t)=\|u_x^\eps(t)\|_\infty$
otrzymujemy na każdym przedziale zaczynającym się
w chwili $r$
\[
 z(t)\le z(r)+CL\eps^{-1/2}
 \int_r^t(t-s)^{-1/2}z(s)\,\mathrm d\lambda_1(s).
\]
Wybór długości przedziału takiej, by
$2CL\sqrt{(t-r)/\eps}\le1/2$, daje
$\sup z(t)\le2z(r)$ na tym przedziale.
Iteracja wyklucza utratę ograniczoności pierwszej
pochodnej w skończonym czasie. Ponieważ norma supremum
rozwiązania jest również ograniczona (z konstrukcji
w $C([0,T];X)$), lokalne twierdzenie o punkcie stałym
można ponawiać z czasem istnienia zależnym tylko od
tych dwóch ograniczeń, $\eps$ i strumienia.
Uzasadnia to kontynuację regularnego rozwiązania
na cały przedział $[0,T]$. Aby uzyskać drugą
pochodną, zapisujemy wzór Duhamela w postaci
z wymuszeniem $F=-\widetilde f'(u^\eps)u_x^\eps$.
Dla $0<\alpha<1$ oszacowania jąder
$\|\partial_xG_\tau\|_1\le C\tau^{-1/2}$ oraz
$\|\partial_{xx}G_\tau\|_1\le C\tau^{-1}$ dają
\[
 \|\partial_xG_\tau(\cdot+h)-\partial_xG_\tau\|_1
 \le C_\alpha |h|^\alpha\tau^{-(1+\alpha)/2}.
\]
Istotnie bierzemy minimum oszacowań przez
$2\|\partial_xG_\tau\|_1$ i
$|h|\|\partial_{xx}G_\tau\|_1$.
Wykładnik $(1+\alpha)/2$ jest mniejszy od $1$,
zatem $u_x^\eps$, a w konsekwencji $F$, są
jednostajnie hölderowskie względem $x$ na skończonym
przedziale czasu (dane Gaussa mają ograniczone
pochodne). Ponieważ $\int_\R\partial_{xx}G_\tau=0$,
przy obliczaniu drugiej pochodnej całki Duhamela
odejmujemy $F(s,x)$ pod całką przestrzenną.
Otrzymujemy całkowalne w czasie oszacowanie
$C_{\eps,\alpha}(t-s)^{-1+\alpha/2}$ dzięki
skalowaniu jądra Gaussa. Aby uzasadnić różniczkowanie, najpierw obcinamy
całkę Duhamela do $0\le s\le t-\eta$, $\eta>0$.
Pochodne rzędu drugiego tych całek zbiegają
jednostajnie (na każdym zwartym przedziale dodatnich
czasów), ponieważ brakujący fragment jest ograniczony
przez $C_{\eps,\alpha}\int_0^\eta
r^{-1+\alpha/2}\,\mathrm d\lambda_1(r)$.
Całki i ich pierwsze pochodne również zbiegają
jednostajnie, z oszacowań jąder rzędu $0$ i $1$.
Zatem granica ma ciągłą drugą pochodną przestrzenną.
Wzór Duhamela spełnia równanie z wymuszeniem $F$
w sensie dystrybucji (testowanie i Fubini, następnie
równanie ciepła dla jądra i jego ślad w chwili $0$).
Prawa strona $H=\eps u_{xx}^\eps+F$ jest ciągła.
Dla ustalonego $r>0$ funkcja
$u^\eps(t,x)-\int_r^t H(s,x)\,\mathrm d\lambda_1(s)$
ma zerową pochodną dystrybucyjną względem $t$.
Testując oddzielnie względem $x$, a następnie
korzystając z ciągłości, wnioskujemy, że jest równa
$u^\eps(r,x)$. Stąd $\partial_t u^\eps=H$
punktowo i rozwiązanie należy do klasy $C^{1,2}$.

Zasada maksimum dla
$u_t^\eps-\eps u_{xx}^\eps+
\widetilde f'(u^\eps)u_x^\eps=0$ daje
$|u^\eps|\le B$. Na całej prostej można ją zastosować
z barierą $\delta e^{At}(1+x^2)$, gdzie
$A>L+2\eps+1$: operator paraboliczny tej bariery
jest ściśle dodatni, bo $2L|x|\le L(1+x^2)$.
Funkcja $u^\eps-B$ pomniejszona o barierę jest
ujemna początkowo i dla dużych $|x|$, jednostajnie
na skończonym przedziale czasu. Gdyby osiągała
dodatnie maksimum na takim walcu, to w punkcie
maksimum (także na jego górnej podstawie) pochodna
czasowa od lewej byłaby nieujemna, pochodna
przestrzenna zerowa, a druga pochodna przestrzenna
niedodatnia. Operator paraboliczny miałby zatem
wartość nieujemną, wbrew jego ścisłej ujemności.
Analogicznie rozpatrujemy $-u^\eps-B$, po czym
przechodzimy z $\delta$ do $0$. Rozwiązanie przyjmuje
więc wartości w przedziale, na którym
$\widetilde f=f$, i rozwiązuje pierwotne równanie.
W szczególności $u^\eps\in C([0,T];L^1(\R))$.

\textbf{Krok 2: kontrakcja lepkościowa.}
Dla dwóch rozwiązań $a,b$ klasy $C^{1,2}$ równania lepkościowego
ich różnica $w=a-b$ spełnia
$w_t+\partial_x(f(a)-f(b))=\eps w_{xx}$.
Przyjmijmy $\eta_\delta(w)=\sqrt{w^2+\delta^2}-\delta$.
Po pomnożeniu przez $\eta_\delta'(w)$ i całkowaniu przez
części, korzystając z $\eta_\delta''\ge0$, otrzymujemy
w granicy $\delta\downarrow0$ nierówność Kato
\[
 \partial_t|a-b|+\partial_x\bigl[
 \operatorname{sgn}(a-b)(f(a)-f(b))\bigr]
 \le\eps\partial_{xx}|a-b|.
\]
Aby uzasadnić przejście do granicy, piszemy
$f(a)-f(b)=A(a,b)(a-b)$, gdzie
$A(a,b)=\int_0^1f'(b+\theta(a-b))\,\mathrm d\theta$.
Dokładna tożsamość ma postać
\[
 \partial_t\eta_\delta(w)+\partial_x(A\eta_\delta(w))
 =\eps\partial_{xx}\eta_\delta(w)
 -\eps\eta_\delta''(w)|w_x|^2
 -\bigl(w\eta_\delta'(w)-\eta_\delta(w)\bigr)\partial_x A.
\]
Ponieważ $w\eta_\delta'(w)-\eta_\delta(w)
=\delta-\delta^2/\sqrt{w^2+\delta^2}\in[0,\delta]$,
ostatni wyraz dąży do zera lokalnie w $L^1$.
Pozostałe wyrazy przechodzą do granicy po testowaniu,
z pominięciem niedodatniego składnika dyssypacyjnego.
Jeśli $a-b\in C([0,T];L^1(\R))$ i oba rozwiązania
są ograniczone (jak rozwiązania z kroku 1), to
$|A|\le M$ dla pewnego $M$. Testujemy otrzymaną
nierówność odcięciem $\chi_{[s,t]}$ wygładzonym
w czasie oraz przestrzenną funkcją $\zeta_R$ równą
$1$ na $[-R,R]$, o nośniku w $[-2R,2R]$,
$0\le\zeta_R\le1$,
$|\zeta_R'|\le C/R$, $|\zeta_R''|\le C/R^2$.
Po granicy czasowej
\[
 \int_\R |w(t,x)|\zeta_R(x)\,\mathrm d\lambda_1(x)
 \le \int_\R |w(s,x)|\zeta_R(x)\,\mathrm d\lambda_1(x)
 +C\left(\frac M R+\frac{\eps}{R^2}\right)
   \int_s^t\|w(\tau)\|_1\,\mathrm d\lambda_1(\tau).
\]
Całka po czasie jest skończona dzięki wcześniejszej
konstrukcji w $C([0,T];L^1)$. Zbieżność
zmajoryzowana dla $R\to\infty$ i ciągłość przy $s=0$
dają $\|a(t)-b(t)\|_1\le\|a(0)-b(0)\|_1$.
Stosując to do $b\equiv0$, otrzymujemy ograniczenie
normy $L^1$. Stosując do przesunięć $u^\eps(t,x+h)$
i $u^\eps(t,x)$, dostajemy
\[
 \|u^\eps(t,\cdot+h)-u^\eps(t,\cdot)\|_1
 \le\|u_0^\eps(\cdot+h)-u_0^\eps\|_1
 \le |h|\operatorname{TV}(u_0).
\]
Aby otrzymać oszacowanie wariacji bez dodatkowego
założenia o całkowalności pochodnej, dla
$\phi\in C_c^1(\R)$, $\|\phi\|_\infty\le1$,
zapisujemy
\[
 \int_\R u^\eps(t,x)\phi'(x)\,\mathrm d\lambda_1(x)
 =\lim_{h\to0}\int_\R
 \frac{u^\eps(t,x-h)-u^\eps(t,x)}{h}\phi(x)
 \,\mathrm d\lambda_1(x)
 \le\operatorname{TV}(u_0).
\]
Równość wynika z podstawienia zmiennej i zbieżności
ilorazów różnicowych $\phi$; oszacowanie z poprzedniej
nierówności. Supremum po $\phi$ daje
$\operatorname{TV}(u^\eps(t))\le\operatorname{TV}(u_0)$.

\textbf{Krok 3: jednostajny moduł ciągłości w czasie.}
Z tożsamości Duhamela stosowanej od czasu $t$ do $t+h$
i z $\|\partial_x f(u^\eps(t))\|_{\mathcal M}\le
M\operatorname{TV}(u_0)$, gdzie
$M=\max_{|z|\le\|u_0\|_\infty}|f'(z)|$, wynika
\[
 \|u^\eps(t+h)-u^\eps(t)\|_1
 \le C\operatorname{TV}(u_0)\sqrt{\eps h}
 +Mh\operatorname{TV}(u_0),\qquad 0<\eps\le1.
\]
Dla $w=u^\eps(t)$ twierdzenie Fubiniego i Lemat
~\ref{lem:pde24-bv-translation} dają
\[
 \begin{aligned}
 \|G_{\eps h}*w-w\|_1
 &\le \int_\R G_{\eps h}(z)\|w(\cdot-z)-w\|_1\,
       \mathrm d\lambda_1(z)\\
 &\le \operatorname{TV}(w)
       \int_\R |z|G_{\eps h}(z)\,\mathrm d\lambda_1(z)\\
 &=2\operatorname{TV}(w)\sqrt{\eps h/\pi}.
 \end{aligned}
\]
Ostatnia równość wynika z podstawienia
$r=z^2/(4\eps h)$ na dodatniej półprostej.
Dla $s>0$ reguła łańcuchowa i ograniczenie
$|f'|\le M$ dają
$\|\partial_x f(u^\eps(s))\|_1
\le M\|u_x^\eps(s)\|_1
=M\operatorname{TV}(u^\eps(s))$;
ostatnia równość wynika z reprezentacji $Dw$ jako
miary ze znakiem. Zatem z nierówności Younga
\[
 \left\|\int_t^{t+h}G_{\eps(t+h-s)}*
 \partial_x f(u^\eps(s))\,\mathrm d\lambda_1(s)\right\|_1
 \le Mh\operatorname{TV}(u_0).
\]
Oba oszacowania razem dają żądany moduł ciągłości.

\textbf{Krok 4: zwartość i granica.}
Uzasadnimy względną zwartość lokalną bez dodatkowego
twierdzenia o zwartości w przestrzeni $BV$.
Dzielimy $[-R,R]$ na $N$ równych przedziałów $I_j$
o długości $\ell=2R/N$ i oznaczamy przez $P_Nw$
funkcję równą średniej $w$ na każdym $I_j$.
Korzystając z nierówności trójkąta, Fubiniego
i Lematu~\ref{lem:pde24-bv-translation}, otrzymujemy
\[
 \begin{aligned}
 \|w-P_Nw\|_{L^1([-R,R])}
 &\le \frac1\ell\sum_{j=1}^N
 \int_{I_j}\!\int_{I_j}|w(x)-w(y)|\,
 \mathrm d\lambda_1(y)\,\mathrm d\lambda_1(x)\\
 &\le\frac1\ell\int_{-\ell}^{\ell}
    \|w(\cdot+h)-w\|_1\,\mathrm d\lambda_1(h)\\
 &\le\ell\operatorname{TV}(w).
 \end{aligned}
\]
W drugiej nierówności podstawiamy $h=y-x$;
dla ustalonego $h$ przedziały całkowania po $x$
są rozłączne i zawarte w $\R$.
Zbiór funkcji $P_Nu^\eps(t)$ jest ograniczony
w skończenie wymiarowej przestrzeni funkcji
stałych na $I_j$ (norma $L^1$ nie wzrasta przy $P_N$).
Dla dowolnej zadanej dokładności najpierw wybieramy
$N$ duże, a potem skończoną sieć w tym ograniczonym
zbiorze. Otrzymujemy całkowitą ograniczoność rodziny
$u^\eps(t)|_{[-R,R]}$, a z zupełności $L^1([-R,R])$
wynika jej względna zwartość. Wspólny moduł ciągłości w
czasie oraz diagonalny argument Arzeli--Ascoliego dają
zbieżność podciągu w $C([0,T];L^1_{\mathrm{loc}}(\R))$.
Aby przejść do zbieżności globalnej, wybieramy
$\zeta_R\in C^\infty(\R)$ równą $0$ na $[-R,R]$ i
równą $1$ poza $[-2R,2R]$, przy
$|\zeta_R'|\le C/R$, $|\zeta_R''|\le C/R^2$.
Nierówność Kato dla $a=u^\eps$, $b=0$ daje
\[
 \int_\R|u^\eps(t,x)|\zeta_R(x)\,\mathrm d\lambda_1(x)
 \le\int_\R|u_0^\eps(x)|\zeta_R(x)\,\mathrm d\lambda_1(x)
 +CT\|u_0\|_1\left(\frac M R+\frac1{R^2}\right).
\]
Testy niegładkie względem czasu i o niezwartym
nośniku względem zmiennej przestrzennej uzasadniamy przez wygładzenie czasu i dodatkowe
odcięcie na $[-N,N]$, po czym $N\to\infty$; pochodne
dodatkowego odcięcia mają rząd $N^{-1}$ i $N^{-2}$,
a całki błędów znikają dzięki oszacowaniu normy $L^1$.
Dla danych początkowych twierdzenie Fubiniego
i nierówność trójkąta dają
\[
 \int_{|x|>R}|u_0^\eps(x)|\,\mathrm d\lambda_1(x)
 \le\int_{|y|>R/2}|u_0(y)|\,\mathrm d\lambda_1(y)
 +\|u_0\|_1\int_{|z|>R/2}G_{\eps^2}(z)\,
      \mathrm d\lambda_1(z).
\]
Ostatnią całkę ograniczamy przez
$4R^{-2}\int_\R z^2G_{\eps^2}(z)\,
\mathrm d\lambda_1(z)=8\eps^2/R^2\le8/R^2$.
Zatem ogony rozwiązań znikają jednostajnie
w $\eps\in(0,1]$ i $t\in[0,T]$. Zbieżny lokalnie
podciąg jest więc ciągiem Cauchy'ego w
$C([0,T];L^1(\R))$: na $[-2R,2R]$ używamy zbieżności
lokalnej, a poza nim sumy dwóch oszacowań ogonów.
Zupełność tej przestrzeni daje zbieżność globalną.

\textbf{Krok 5: równanie i nierówności entropijne.}
Wzór słaby dla $u^\eps$ przechodzi do granicy, gdyż
$f$ jest lipschitzowska na wspólnym zakresie wartości,
a składnik lepkościowy
\[
 \eps\int_0^T\!\int_\R u^\eps\varphi_{xx}\,
 \mathrm d\lambda_1(x)\,\mathrm d\lambda_1(t)
\]
dąży do zera. Dla gładkiej wypukłej entropii
korzystamy z~\eqref{eq:pde24-entropyvisc}.
Po przetestowaniu nieujemną $\varphi$ otrzymujemy
\[
 \int_0^T\!\int_\R
 \bigl(\eta(u^\eps)\varphi_t+q(u^\eps)\varphi_x+
       \eps\eta(u^\eps)\varphi_{xx}\bigr)
 \,\mathrm d\lambda_1(x)\,\mathrm d\lambda_1(t)\ge0.
\]
Ostatni składnik dąży do zera dzięki wspólnemu
ograniczeniu $u^\eps$ i zwartemu nośnikowi testu.
Funkcje $\eta,q$ są lipschitzowskie na wspólnym
przedziale wartości, więc silna zbieżność w
$L^1_{\mathrm{loc}}$ pozwala przejść do granicy
w pozostałych składnikach. Lemat
~\ref{lem:pde24-entropy-equivalence} daje warunki
Krużkowa. Zbieżność w $C([0,T];L^1)$ oraz
$G_{\eps^2}*u_0\to u_0$ w $L^1$ gwarantują $u(0)=u_0$. Ograniczenia~\eqref{eq:pde24-BVbound}
przechodzą do granicy w każdej chwili $t$:
norma $L^1$ jest ciągła, dla ograniczenia $L^\infty$
wybieramy podciąg zbieżny prawie wszędzie w tej
ustalonej chwili, a wariacja jest dolnie półciągła,
gdyż jest supremum ciągłych funkcjonałów
$w\mapsto\int_\R w\phi'\,\mathrm d\lambda_1$.
Jednoznaczność wynika z Twierdzenia~\ref{thm:pde24-L1}.
Zbiega cała rodzina: w przeciwnym razie istniałby
ciąg $\eps_n\downarrow0$ o odległości od $u$
w $C([0,T];L^1)$ ograniczonej od dołu przez dodatnią
stałą. Poprzedni argument zwartości dałby podciąg
zbieżny do rozwiązania entropijnego, które z
jednoznaczności musi być równe $u$, sprzeczność.
\end{proof}


\begin{twierdzenie}[Istnienie dla danych $L^1\cap L^\infty$ bez założenia $BV$]
\label{thm:pde24-existence-L1Linf}
Niech $f\in C^2(\R)$, $u_0\in L^1(\R)\cap L^\infty(\R)$
i $T>0$. Wtedy istnieje dokładnie jedno rozwiązanie entropijne
$u\in C([0,T];L^1(\R))\cap L^\infty((0,T)\times\R)$
z warunkiem $u(0)=u_0$. Spełnia ono
\[
 \|u(t)\|_\infty\le\|u_0\|_\infty,\qquad
 \|u(t)\|_1\le\|u_0\|_1
 \quad(0\le t\le T).
\]
Jeżeli dodatkowo $u_0\in BV(\R)$, to
$\operatorname{TV}(u(t))\le\operatorname{TV}(u_0)$.
\end{twierdzenie}
\begin{proof}
Wybieramy ciąg $u_0^n\in C_c^\infty(\R)$ taki, że
$u_0^n\to u_0$ w $L^1(\R)$ oraz
$\|u_0^n\|_\infty\le\|u_0\|_\infty$: wystarczy
obciąć $u_0$ przez funkcje z $[0,1]$, a następnie
zastosować splot z nieujemnym jądrem aproksymującym jedność.
Każde $u_0^n$ należy do $BV\cap L^1\cap L^\infty$,
a z Twierdzenia~\ref{thm:pde24-existence} istnieje
odpowiadające mu rozwiązanie entropijne $u^n$.
Oszacowanie~\eqref{eq:pde24-L1} daje
\[
 \sup_{0\le t\le T}\|u^n(t)-u^m(t)\|_1
 \le\|u_0^n-u_0^m\|_1.
\]
Zupełność $C([0,T];L^1(\R))$ zapewnia więc istnienie
$u$ z $u^n\to u$ w tej przestrzeni. Dla każdej chwili $t$ zbieżność w $L^1$ daje
\[
 \int_\R (|u(t,x)|-B)_+\,\mathrm d\lambda_1(x)
 \le\|u(t)-u^n(t)\|_1\longrightarrow0,
 \qquad B=\|u_0\|_\infty,
\]
gdzie $r_+=\max(r,0)$. Zatem $\|u(t)\|_\infty\le B$
dla każdego $t$, a z Fubiniego otrzymujemy
$\|u\|_{L^\infty((0,T)\times\R)}\le B$.
Ponieważ $f$ i strumienie Krużkowa
$q_k(z)=\operatorname{sgn}(z-k)(f(z)-f(k))$
są lipschitzowskie na wspólnym ograniczonym zakresie
wartości (dla każdego ustalonego $k$), możemy przejść
do granicy we wszystkich nierównościach entropijnych
na testach o zwartym nośniku. Ciągłość w $L^1$
zapewnia ślad $u(0)=u_0$. Ograniczenie normy $L^1$
przechodzi do granicy; jednoznaczność jest konsekwencją
Twierdzenia~\ref{thm:pde24-L1}. Jeżeli $u_0\in BV$,
możemy też skorzystać wprost z Twierdzenia
\ref{thm:pde24-existence}, które daje ograniczenie wariacji.
\end{proof}


\begin{twierdzenie}[Istnienie dla danych ograniczonych]
\label{thm:pde24-existence-Linf}
Niech $f\in C^2(\R)$, $u_0\in L^\infty(\R)$ oraz $T>0$.
Istnieje dokładnie jedna funkcja
\[
 u\in C([0,T];L^1_{\mathrm{loc}}(\R))
 \cap L^\infty((0,T)\times\R)
\]
spełniająca lokalne nierówności Krużkowa
~\eqref{eq:pde24-kruzineq} oraz $u(0)=u_0$
w $L^1_{\mathrm{loc}}(\R)$. Ponadto
$\|u(t)\|_\infty\le\|u_0\|_\infty$ dla każdego $t$.
\end{twierdzenie}
\begin{proof}
Niech $B=\|u_0\|_\infty$,
$M=\max_{|z|\le B}|f'(z)|$ i
$u_0^n=u_0\chi_{[-n,n]}$ dla $n\ge1$.
Z Twierdzenia~\ref{thm:pde24-existence-L1Linf}
otrzymujemy rozwiązania $u^n$ o globalnym
ograniczeniu przez $B$.
Ustalmy $R>0$ i wybierzmy $n,m>R+MT$.
Dla każdego $t\in[0,T]$ Wniosek~\ref{cor:pde24-finite}
daje
\[
 \int_{-R}^R |u^n(t,x)-u^m(t,x)|\,\mathrm d\lambda_1(x)
 \le\int_{-R-Mt}^{R+Mt}|u_0^n(x)-u_0^m(x)|
 \,\mathrm d\lambda_1(x)=0.
\]
A zatem ograniczenia $u^n$ do $[-R,R]$ stabilizują się
w $C([0,T];L^1([-R,R]))$. Dla $R=1,2,\ldots$
otrzymane granice są zgodne przy ograniczaniu
do mniejszych przedziałów. Sklejamy je w funkcję
$u\in C([0,T];L^1_{\mathrm{loc}}(\R))$.
Możemy wybrać wspólnie mierzalnego reprezentanta,
sklejając reprezentanty $u^n$ na kolejnych
pierścieniach przestrzennych; ograniczenie przez $B$
przechodzi na każdą chwilę oraz prawie wszędzie
w czasoprzestrzeni. Każdy test entropijny ma nośnik
w pewnym walcu $(0,T)\times[-R,R]$, gdzie $u=u^n$
prawie wszędzie dla dostatecznie dużego $n$.
Stąd $u$ spełnia wszystkie lokalne nierówności.
Zgodność w chwili $0$ na każdym takim przedziale
daje ślad początkowy. Dwa dowolne rozwiązania
z podanej klasy mają wspólne globalne ograniczenie
(a więc wspólną skończoną prędkość w oszacowaniu
stożkowym), toteż Wniosek~\ref{cor:pde24-finite}
zapewnia ich jednoznaczność.
\end{proof}

\section{Związek z równaniem Hamiltona--Jacobiego i dalsze problemy}
Dla $f(w)=w^2/2$ i funkcji $g\in W^{1,\infty}(\R)$
rozwiązanie lepkościowe równania
$U_t+\tfrac12(U_x)^2=0$ z danymi $U(0,\cdot)=g$
ma pochodną
$w=U_x$, która jest rozwiązaniem entropijnym
równania Burgersa, w odpowiednim sensie dystrybucyjnym.
Dla kwadratowego hamiltonianu jest to już udowodnione
w rozdziale~\ref{chap:pde23}, w sekcji o związku
z równaniem Burgersa i warunkiem entropii, metodą
półwklęsłości rozwiązania Hopfa--Laxa. Otrzymujemy
tam także $w\in C([0,T];L^1_{\mathrm{loc}}(\R))$,
$|w|\le\|g'\|_\infty$ i ślad $w(0)=g'$.
Wniosek~\ref{cor:pde24-finite} zapewnia jednoznaczność
w klasie globalnie ograniczonych rozwiązań lokalnych.
Jeżeli dodatkowo $g'\in L^1(\R)$, to twierdzenie
~\ref{thm:pde24-existence-L1Linf} oraz ta sama lokalna
jednoznaczność identyfikują $w$ z rozwiązaniem
w $C([0,T];L^1(\R))$. Dla dowolnej globalnie
lipschitzowskiej funkcji $g$ (niekoniecznie ograniczonej)
istnienie i jednoznaczność lokalnego rozwiązania
entropijnego wynikają też z Twierdzenia
~\ref{thm:pde24-existence-Linf}, stosowanego do $g'$.
Nie utożsamiamy natomiast
automatycznie granic pochodnych przybliżeń
lepkościowych na podstawie samej zbieżności potencjałów.

\begin{przyklad}[Dyssypacja energii w fali Burgersa]
Dla $f(w)=w^2/2$ wybieramy $\eta(w)=w^2/2$, wtedy
$q(w)=w^3/3$. Przy skoku z $u_L=1$ do $u_R=0$
prędkość wynosi $\sigma=1/2$, a współczynnik
entropijnej delty jest równy
$[q]-\sigma[\eta]=-1/3+1/4=-1/12<0$.
Dla odwróconego skoku współczynnik jest dodatni,
więc warunek entropii zostaje naruszony. Na mocy
Lematu~\ref{lem:pde24-entropy-equivalence} skok ten
nie spełnia również wszystkich nierówności Krużkowa.
\end{przyklad}

\begin{uwaga}[Zakres twierdzeń]
Globalna kontrakcja $L^1$ dotyczy rozwiązań
o całkowalnych i ograniczonych danych; dla oszacowania
wariacji dodatkowo zakładamy $BV$. Twierdzenie
~\ref{thm:pde24-existence-Linf} obejmuje także dane
jedynie ograniczone, w tym dane Riemanna. W tej klasie
stosujemy oszacowanie w stożku zamiast globalnej normy
$L^1$. Wyniki dla układów praw zachowania,
strumieni zależnych od $x$, wyższych wymiarów lub
nie-wypukłych zadań Riemanna wymagają dodatkowych metod.
\end{uwaga}



% ROZDZIAL 25: nowe rownania pololiniowe
